\section{Preliminaries}\label{sec:preliminaries}

This section fixes the mathematical setting and notation used throughout the paper. The central point is that the theorem object of the paper is not the raw simulation trajectory of the continuous substrate, but a canonical object built from two coupled levels of description: a cocycle-pressure theory \(\Tzero\) and a completion-based extension \(\Tone\). Standard background references for the surrounding concepts include classical IFS/fractal geometry \cite{Hutchinson1981Fractals,FalconerFractalGeometry,MauldinUrbanskiGDMS,Moran1946AdditiveFunctions}, thermodynamic formalism \cite{Ruelle1978ThermodynamicFormalism,WaltersErgodic1982,BarreiraNonadditiveTF}, and subadditive or almost-additive extensions of that formalism \cite{CaoFengHuang2008SubAdditive,Mummert2006AlmostAdditive}.

\subsection{Continuous substrate states and the audited shell}

A \emph{continuous substrate state} is written
\[
x=(K,\ell,\mathfrak{p},\tau,b,\eta)\in \mathcal{X},
\]
where:
\begin{itemize}
    \item \(K\in\mathcal{K}\subset \R^{d\times d}\) is the current normalized kernel/operator state,
    \item \(\ell\in\mathcal{L}\) is the currently active lens state,
    \item \(\mathfrak{p}\in\mathcal{P}\) is the currently active packaging state,
    \item \(\tau\in\mathcal{T}\subset \R_{>0}\) is the current timescale/protocol parameter,
    \item \(b\in\mathcal{B}\subset \R\) is the budget/ledger variable,
    \item \(\eta\in\mathcal{H}\) denotes the remaining auxiliary selector and hysteresis state.
\end{itemize}
The full update law is written
\[
F:\mathcal{X}\to\mathcal{X},\qquad x\mapsto F(x),
\]
and is generated by the six interacting mechanisms listed in the introduction:
operator rewrite (\(P1\)), admissibility gating (\(P2\)), protocol/timescale adaptation (\(P3\)), lens selection (\(P4\)), packaging/completion (\(P5\)), and budget/ledger dynamics (\(P6\)).

We work on a distinguished audited shell-stable class
\[
\mathcal{S}_{\mathrm{aud}}\subset \mathcal{X},
\]
which is the class denoted in the internal package by
\[
\texttt{continuous\_full\_loop\_lawful\_kernel\_class\_shell\allowbreak\_stable}.
\]
All theoremlets in this paper are statements on \(\mathcal{S}_{\mathrm{aud}}\).

\subsection{The base theory \texorpdfstring{$\Tzero$}{T0}}

The base theory \(\Tzero\) is the cocycle-pressure theory carried by the audited shell. We write
\[
\pi_0:\mathcal{S}_{\mathrm{aud}}\to \mathcal{O}_0
\]
for the \(\Tzero\)-object map, where \(\mathcal{O}_0\) denotes the cocycle-level object space. Intuitively, \(\pi_0(x)\) records the object visible at the cocycle/lens level before completion-based packaging is adjoined.

The thermodynamic object attached to \(\Tzero\) is a cocycle-pressure quantity. For the chosen observable family \(a_n(s;x)\), depending on depth/time \(n\), parameter \(s\), and shell state \(x\), the associated pressure is written
\[
P_{\Tzero}(s)
\]
once the closure theoremlet has established that the corresponding pressure object is well-defined on the audited shell. At the level of notation we suppress the shell-state dependence whenever the shell-stable theorem package has already fixed it.

\subsection{The completion extension \texorpdfstring{$\Tone$}{T1}}

The extended theory \(\Tone\) adjoins packaging completion to the cocycle theory. For a given lens \(\ell\) and timescale \(\tau\), the packaging endomap is
\[
E_{\tau,\ell}(\mu):=U_{\ell}\!\bigl(Q_{\ell}(\mu K^{\tau})\bigr),
\]
where:
\begin{itemize}
    \item \(\mu\) is a state/distribution object on which the completion acts,
    \item \(Q_{\ell}\) is the lens-dependent forgetting/coarse-graining map,
    \item \(U_{\ell}\) is the corresponding reinstatement/prototype lift.
\end{itemize}
The fixed points of \(E_{\tau,\ell}\) are the packaged objects of the completion theory. We write
\[
\mathrm{Fix}(E_{\tau,\ell})
\]
for the fixed-point set and
\[
\pi_1:\mathcal{S}_{\mathrm{aud}}\to \mathcal{O}_1
\]
for the \(\Tone\)-object map, where \(\mathcal{O}_1\) is the packaged-object space.

A \emph{packaged stratum} is a persistent element or equivalence class in the packaged fixed-point structure associated with the current shell state. We write
\[
\mathcal{F}(x)
\]
for the finite family of persistent packaged strata determined by \(x\in \mathcal{S}_{\mathrm{aud}}\).

\subsection{The canonical hybrid object}

The theorem object of the paper is not \(\pi_0(x)\) alone and not \(\pi_1(x)\) alone. It is the canonical hybrid object obtained by combining the cocycle-level object map, the packaging-completion endomap, and the packaged strata. We write it schematically as
\[
\mathfrak{H}(x)=\bigl(\pi_0(x),\,\pi_1(x),\,\mathcal{F}(x)\bigr).
\]
The role of \(\mathfrak{H}\) is twofold:
\begin{enumerate}
    \item it keeps track of the closed cocycle-pressure theory \(\Tzero\), and
    \item it records the genuinely richer packaged-object structure used to define the extension \(\Tone\).
\end{enumerate}
In particular, strict theory extension will be formulated as a statement that the packaged-object map does not factor through the cocycle-level object map on the audited shell.

\subsection{Pressure and conditional disintegration objects}

The cocycle pressure \(P_{\Tzero}(s)\) is the base thermodynamic quantity. On the completion side, the paper does \emph{not} use direct stratumwise root separation as its thermodynamic consequence. Instead, the consequence object is a weighted package-conditioned pressure gap. If \(\mathcal{F}(x)\) is the family of packaged strata and \(w_{\Sigma}(x)\) are the associated package weights, then the consequence quantity is written
\[
\Delta_{\Tzero\to \Tone}(s)
:=
P_{\Tzero}(s)-\sum_{\Sigma\in\mathcal{F}(x)} w_{\Sigma}(x)\,P_{\Sigma}(s),
\]
where \(P_{\Sigma}(s)\) denotes the package-conditioned pressure profile attached to the stratum \(\Sigma\). The conditional-disintegration theoremlet states that this weighted gap is nontrivial on the audited shell.

The role of \(\Delta_{\Tzero\to \Tone}(s)\) is that it measures the thermodynamic insufficiency of the cocycle-level theory for the packaged future. In the paper’s final consequence theoremlet, the relevant thermodynamic statement is therefore a \emph{conditional disintegration statement}, not a direct theorem of independent stratumwise root separation \cite{Rokhlin1949FundamentalIdeas}.

\subsection{Definability, saturation, and obstruction}

Two further pieces of notation are used repeatedly.

First, \(\Tzero\)-\emph{definability} means definability from the cocycle-level object/lens structure alone. At the formal level this is expressed through factorization of the packaged-object map through \(\pi_0\). Failure of such factorization on the audited shell is the non-definability signal behind the strict theory extension theoremlet.

Second, the completion dynamics may \emph{saturate}: repeated application of \(E_{\tau,\ell}\) ceases to produce new packaged objects at the current theory depth. Material \(P4\leftarrow P5\) forcing then changes the lens and reveals new packaged strata in the extended theory.

Finally, we use \emph{macro-admissibility obstruction} to mean that the cocycle-level theory is too coarse to support closed macro-dynamics on the packaged future. This obstruction is one of the certificates that \(\Tone\) is a strict theory extension of \(\Tzero\).

\subsection{Scope convention}

Unless stated otherwise, every theorem in the paper is made on the audited shell \(\mathcal{S}_{\mathrm{aud}}\) and for the canonical hybrid object \(\mathfrak{H}\). The explicit list of non-claims is deferred to Section~\ref{sec:discussion}.
